\section{Evaluation}
\label{sec:eval}
% Data: baseline checkpoint run_20260616_011948; 16 archived KernelSeek sessions
% (paper/analysis/traces); figures from scripts/plot_eval.py.

We ask five questions, one per hypothesis of \S\ref{sec:hypotheses}. RQ1: how fast and how correct
are the kernels? RQ2: on the same task, evaluator, hardware, and week, how does \sys compare with the
equipped agent? RQ3: which human messages redirected the agent, and how? RQ4: what does evolutionary
search add on top of questions? RQ5: do the questions transfer?

\subsection{Setup}
\label{sec:eval-setup}

Hopper experiments and the head-to-head ran on one 8$\times$H200 node; Blackwell experiments on one
8$\times$B200 node. The upstream kernels are vLLM's FA3 and the CuteDSL FA4 at pinned commits;
cuDNN's fused attention (9.20 on Hopper, 9.19 on Blackwell) and FlashInfer are comparison columns
fed the same buffers. Implementations are timed back to back on one idle GPU with an L2 flush before
each sample; clocks are not locked, so only ratios are reported; four independent Hopper batches
agree to within 0.3\% on every aggregate (Appendix~\ref{app:bench}). The equipped
agent ran \texttt{claude-opus-4-8} through Claude Code in headless mode, one fresh session per step.
\sys sessions used two front-ends: Claude Code with the Stop-hook goal, on \texttt{claude-opus-4-8}
for seven sessions and \texttt{claude-opus-5} for the last two, and Cursor with the same instructions
and without the hook, whose exports do not record the model. Sixteen \sys sessions are archived, nine
in Claude Code and seven in Cursor, from June 17 to August 7, 2026: twelve for the Hopper campaign
and four for Blackwell. They contain 162 human messages and 11,532 agent tool calls, 8,791 in the
parent sessions and 2,741 in subagents. One author labeled each human message by its primary act
(goal, directive, question, approval, correction), flagged interrogative phrasing, and assigned
questions to the eight categories of \S\ref{sec:workflow-questions}; tool calls were classified by
target (Appendix~\ref{app:traces}).

\subsection{RQ1: Kernel Performance}
\label{sec:eval-perf}

The kernels are reported in \S\ref{sec:kernels}. In brief: \fathree is correct on all 182 Hopper
scenarios at 1.16$\times$ FA3 and 1.10$\times$ the fastest available implementation including cuDNN;
\fafour is correct on all 154 runnable Blackwell scenarios at 1.46$\times$ FA4 and 1.07$\times$ the
fastest available. Per-scenario Hopper ratios vary by 1\% at the median and 6\% at most across
repeats. Every Blackwell scenario below parity has a measured root cause (Appendix~\ref{app:kernels}).
This is H1's quality half, and the result the rest of the section explains.

\subsection{RQ2: Head-to-Head on the Same Task}
\label{sec:eval-h2h}

\begin{table}[t]
\centering\small
\caption{Head-to-head on the equipped agent's own evaluator: 26 FA3 configurations, same H200 node,
overlapping days in June 2026. Upper block: outcome; the \sys hours run from the message that set
the 26-configuration goal to the best result, excluding a 26.7-hour idle gap. Lower block: tool use;
equipped-agent figures are over its 133 steps, \sys figures over all 16 sessions.}
\label{tab:h2h}
\setlength{\tabcolsep}{3pt}\footnotesize\raggedright
\begin{tabular}{@{}p{0.40\columnwidth}>{\raggedright\arraybackslash}p{0.25\columnwidth}>{\raggedright\arraybackslash}p{0.28\columnwidth}@{}}
\toprule
 & Equipped agent & \sys \\
\midrule
Human text in prompts       & none          & 40 messages, 17 questions \\
Model-driven steps / turns  & 133 steps     & 40 turns \\
Wall-clock to best          & 44.7 h of 62.4 h & 68.0 h, 41.3 h active \\
API cost                    & \$660         & not recoverable (Cursor) \\
Candidates correct          & 132 / 133     & all reported \\
Best geomean                & 0.48$\times$  & 1.13$\times$ \\
Worst configuration         & 0.24$\times$  & 0.87$\times$ \\
Configurations $\ge$1.0$\times$ & 0 / 26    & 22 / 26 \\
Goal met                    & no            & yes \\
\midrule
Reference material pushed into prompt & 37\% of prompt & none \\
Tool-result text re-reading own kernel & 82\% & 26\% \\
Profiler (NCU) runs         & 1 per step, parsed & 429, raw \\
SASS / PTX inspections      & 1 per step, parsed & 372 \\
Documentation reads         & 92 lookups    & 36 \\
Builds / benchmark runs     & 133 / 133     & 377 / 937 \\
Subagents spawned           & 0             & 48 \\
\bottomrule
\end{tabular}
\end{table}

The comparison is controlled in the ways that matter. During its first Hopper session, \sys adopted
the equipped agent's own evaluator, the same 26 configurations, FA3 baseline, and correctness gate,
on the same machine, in overlapping days. Table~\ref{tab:h2h} and Figure~\ref{fig:teaser} give the
outcome. The equipped agent's lineage champion reached 0.48$\times$ FA3 at step 35, after 44.7
hours, and did not improve in the remaining 98 model-driven steps (\S\ref{sec:baseline}). Its worst
configuration stayed at 0.24$\times$ and none reached parity. \sys started from 0.009$\times$,
because its kernel of the time routed the evaluator's paged and grouped-query configurations to the
scalar fallback, and reached 0.96$\times$ after 21 hours, 1.02$\times$ after 66, and 1.13$\times$
after 68, of which 26.7 were an idle gap. At the end, 22 of 26 configurations were at or above parity
and the worst was 0.87$\times$. This supports H1 and, counting the 40 human turns against 133 steps
and 62 model-hours, H2.

Three caveats qualify the comparison: \sys had built a kernel for FA3's 40-scenario test grid in
the two preceding days, so the 0.009$\times$ start understates its head start; the front-ends
differ, and token usage is not recoverable for the Cursor sessions; and each arm is a single run.
The fair reading is that in one calendar week a steered agent crossed parity on an evaluator on
which an unsteered one, given more compute, did not reach half of it. We treat Table~\ref{tab:h2h}
as showing that the gap exists and is large, and RQ3 as the evidence for why.

\paragraph{Where the effort went.} The lower block of Table~\ref{tab:h2h} contrasts tool use. The
equipped agent received a third of its context as pre-fetched reference material, yet spent 82\% of
what it read on its own kernel. \sys, with nothing pushed to it, spent 26\% on its own kernels and
ran the profiler 429 times on raw output. The questions did not reduce tool use; they aimed it.

\subsection{RQ3: Anatomy of the Questions}
\label{sec:eval-cases}

\paragraph{What the human said.} Of the 162 human messages, directives are the largest class (58,
36\%), followed by questions (44, 27\%), approvals (30, 19\%), corrections (19, 12\%), and goals
(11, 7\%); 63 messages (39\%) are interrogative in phrasing, and 12 of the 19 corrections are
phrased as questions (Appendix~\ref{app:traces}). The steering was interrogative, not mostly questions:
directives and approvals keep the agent moving, and the questions are where its course changed. The
agent also asked. Across the Claude Code sessions it
posed 35 structured questions to the human, 24 of which offered an option to accept the current
state or stop, and its messages assert a limit or an impossibility 126 times. The Stop hook returned
it to work 59 times. Autonomous agents seek closure; the harness and the human refused it.

\begin{table*}[t]
\centering\small
\caption{Case studies from the archived traces: the question (paraphrased for length), what the
agent believed before it, the measurement the question triggered, and the outcome. Speedups are
vs.\ the upstream kernel on the evaluator in use at the time.}
\label{tab:cases}
\setlength{\tabcolsep}{4pt}\footnotesize
\begin{tabular}{@{}p{0.22\textwidth}p{0.22\textwidth}p{0.24\textwidth}p{0.26\textwidth}@{}}
\toprule
Question (category) & Agent's belief & Measurement it triggered & Outcome \\
\midrule
``What configs are bringing down the geomean?'' (gap triage) &
Tuning the hd128 paths was the right priority. &
Per-configuration table: hd96, hd192, hd256 were on the scalar fallback at 0.01$\times$. &
Fallback removed; 0.44 $\to$ 0.96 geomean in 4 hours. \\
\addlinespace
``Is this ceiling a property of the hardware, or of our design? Where does the time go outside the kernel?'' (bound differencing) &
Split-KV decode was at a ``genuine ceiling within the harness rules''. &
Host-side profile: a per-call tensor fill of 2--5~$\mu$s outside the kernel. &
Fill removed; split-KV 0.87 $\to$ 0.97, geomean 1.02 $\to$ 1.13. \\
\addlinespace
``Do our scenarios cover the shapes and features the operator is actually called with?'' (benchmark validity) &
40 of 40 scenarios correct meant the kernel was done. &
The operator's full test grid ported: 52 of 58 correct, with sink and left-pad outputs off by 495$\times$ and 103$\times$. &
Six fixes; 58 of 58 correct, geomean 1.10 on the new grid. \\
\addlinespace
``Why do you think the goal was achieved? We are still well below parity.'' (challenge a limit) &
The Stop-hook judge had accepted ``goal met''. &
Agent: ``I stopped after the first successful lever and declared victory prematurely.'' &
fp8 hd256 decode 0.20/0.34 $\to$ 0.66/0.77 after ``this is not a rare corner''. \\
\addlinespace
``Is it NCU that hangs, or our kernel?'' (hypothesis check) &
``No deadlock in our kernel.'' &
Sanitizer, then debugger: a warp-convergent \texttt{tcgen05.wait} inside a divergent branch. &
All 8 hanging shapes fixed. \\
\addlinespace
``Is this really a register floor?'' (challenge a limit) &
hd256 was register-capped; the judge had accepted the limit twice. &
A collective rewrite of the consumer loop. &
0.84 $\to$ 0.95--0.96$\times$; the agent's own note: ``the `floor' conclusion was wrong''. \\
\bottomrule
\end{tabular}
\end{table*}

\paragraph{Which questions mattered.} Table~\ref{tab:cases} walks through six turning points we can
verify in the archived traces. Two patterns stand out. First, every large jump followed a question
that either moved the agent to a different layer (from kernel tuning to the per-configuration table,
to the host launch path, to the benchmark itself) or rejected a claimed limit; none followed a
message that supplied information the agent lacked, which supports H3. Second, the measurement the
question triggered was always one the agent could already make: the per-configuration table, the
host-side profile, the sanitizer, and the debugger were available at every earlier step. What the
agent lacked was a reason to look. Four further episodes, among them the 38~$\mu$s host copy and a
routing sweep that moved the Blackwell geomean from 0.88 to 1.15$\times$ without changing a kernel
body, are documented in the repository's history but not in archived sessions, and we cite them as
outcomes only.

\subsection{RQ4: Questions and Evolutionary Search}
\label{sec:eval-evo}


We wrapped the same agent in the loop of \S\ref{sec:workflow-evo} and ran it on the Blackwell
dense-prefill kernel from a correct but slow port (Figure~\ref{fig:evo}, Appendix~\ref{app:traces}). The first generation,
whose directive asked whether a persistent schedule with longest-first ordering and cross-tile
overlap would recover the lost overlap, took the kernel from 0.29$\times$ to 0.68$\times$. Three
more generations added 0.10$\times$ between them through a folded softmax scale, a split K/V
barrier, and a depth-two softmax--correction handshake, while directives on register rebalancing and
address hoisting measured neutral or negative. The search stalled at 0.78$\times$ with the kernel's
own cycle count already below upstream's. The last 22 points came
from a question, not a generation: asked where the time was going, the agent profiled the whole call
rather than the kernel and found that synchronous host copies accounted for the gap. Evolution
amplified the direction the directives carried and searched efficiently within the layer it was
pointed at; it could not move to a layer it was not pointed at. This is H4.

\subsection{RQ5: Reuse}
\label{sec:eval-auto}

The question categories were populated mostly during the Hopper campaign, which came first. Five of
the eight recur in the Blackwell campaign, and the Blackwell skill's self-questions restate Hopper
questions nearly verbatim. Of the 63 interrogative messages, 17 mention a concept specific to
attention and 32 mention the upstream kernel; the rest name only measurements, layers, and process,
and where a question names the upstream it uses it as a bound, a role a roofline fills when no
upstream exists. This is the evidence for H5 that we have: necessary, since a bank that did not
transfer across two architectures would not transfer further, but not sufficient, since both
campaigns were attention and the 72 bank entries mined from public sources (\S\ref{sec:workflow-bank})
have not yet steered an agent. The experiment that would test both, and whether the bank can replace
the human at decision points, is described in \S\ref{sec:discussion}.
